\section{引言}
\label{sec:introduction}

\subsection{理论发展}

赤道波的现代理论框架始于 \citep{matsunoQuasiGeostrophicMotionsEquatorial1966} 的开创性工作。
Matsuno 在赤道 $\rossby$ 平面近似下，对线性化的浅水方程组进行了求解。
他发现，由于科氏参数 $\coriolis$ 随纬度的线性变化，波动解在经向方向上被束缚在赤道附近，其结构由抛物柱函数（厄米函数）描述。
基于这一数学框架，Matsuno 成功推导出了赤道波动的完整频散关系，建立了一套正交完备的波动模态集：包括对应于 $\modenum = -1$ 的 Kelvin 波，对应于 $\modenum = 0$ 的混合 Rossby-重力波，以及对应于 $\modenum \ge 1$ 的赤道 Rossby 波和惯性重力波。
这一理论不仅揭示了赤道地区独特的波动动力学特征，更为此后半个多世纪的热带气象学研究奠定了坚实的基石。

Matsuno 的理论虽然最初是基于干大气的理想化假设推导而来的，但其有效性很快在随后的观测中得到了令人惊叹的验证。
在 20 世纪 60 年代末，利用太平洋地区的无线电探空网络，\citep{Yanai1966} 首先在热带平流层下部识别出了周期约为 4--5 天的 混合 Rossby-重力波。
紧随其后，\citep{Wallace1968} 也在平流层确证了 Kelvin 波的存在，观测显示其表现为单纯的纬向风与位势高度场的同相扰动，且经向风分量几乎为零，这与理论预测完全吻合。
这两种波动被证实具有显著的垂直能量传播特征，能够将对流层的动量向上传输，为后来解释热带平流层准两年振荡（QBO）的驱动机制提供了关键的观测证据。

进入 70 年代，观测的触角开始延伸至更加复杂的热带对流层。
\citep{Madden1979} 利用长时间序列的全球探空资料，首次在对流层低层成功分离出了大尺度的赤道 Rossby 波，其向西传播的相速度和低频特征进一步丰富了赤道波动的观测图谱。
随着卫星遥感时代的到来，研究手段发生了质的飞跃。
\citep{Takayabu1994}利用高分辨率的地球同步卫星红外资料（GMS-IR），率先在波数-频率域中发现热带云团的组织结构并非杂乱无章，而是表现出明显的波动特征。
随后，\citep{Wheeler1999} 将这一谱分析方法推广至全球，通过对多年长波辐射（OLR）资料的处理，剔除了红噪声背景，绘制了著名的波数-频率能量谱图。
他们发现，热带对流活动的显著谱峰精确地沿着 Matsuno 理论导出的线性频散曲线分布，这一结果在统计学上给出了铁证：即使在充满强非线性对流加热的热带对流层中，Matsuno 描绘的线性动力学骨架依然主导着大尺度对流系统的时空演变。
此外，基于 TOGA COARE 等大型科学试验的高空加密观测也证实 \citep{kiladisConvectivelyCoupledEquatorial2009}，这些波动在垂直结构上表现出类似于第一斜压模态的倾斜特征，进一步验证了理论的普适性。

然而，观测同时揭示了经典干波理论的不足。
对流耦合赤道波（CCEWs）的传播速度系统性地低于干波理论的预测，对应的等效深度约为 $12$--$50$ m，远小于干大气的第一斜压模态（$\sim 200$ m）。
这一差异表明，热带对流中的非绝热加热 $\nonadheat$ 对波动结构具有根本性影响。
针对这一挑战，\citep{hayashiTheoryLargeScaleEquatorial1970} 突破了浅水方程的限制，发展了层结大气中赤道波动的三维理论，推导出了包含垂直切变和非绝热加热垂直分布的更一般化的波动方程。
在此基础上，\citep{lindzenWaveCISKTropics1974} 提出了著名的 Wave-CISK（第二类条件不稳定） 机制，为理解湿过程与波动的耦合提供了核心理论框架。
该理论指出，对流潜热释放并非独立过程，而是与大尺度波动的低层辐合形成正反馈：波动辐合触发对流，对流释放潜热进一步加热大气，这种非绝热加热 $\nonadheat$ 改变了波动的有效静力稳定度（Effective Static Stability）。
当加热与波动垂直运动相位适当时，虽然大气的总静力稳定度保持为正，但波动感受到的“有效”稳定度显著降低。
这种热力强迫作用直接导致波动相速度减慢，使其在观测上表现出极浅的等效深度。
这一系列理论进展成功将赤道波动研究从单纯的干流体动力学推进到了“对流耦合动力学”。

\subsection{科学问题与研究动机}

尽管现有理论已取得显著进展，但仍存在若干关键问题。

\begin{enumerate}
    \item 如何从三维层结大气方程严格推导出等效于 Matsuno 的二维浅水系统？等效深度 $\eqdepth$ 作为连接三维与二维理论的核心参数，其数学起源是什么？

    \item 非绝热加热 $\nonadheat$ 的参数化形式如何影响垂直-水平结构的可分性？湿波与干波的本质区别是什么？

    \item 能否建立一个算子代数框架，将不同类型的赤道波动纳入同一数学结构，并利用微扰方法系统分析湿过程的影响？
\end{enumerate}

本文旨在构建一个统一的数学理论，其核心思路是：首先在干大气假设下严格推导垂直-水平分离，识别等效深度作为垂直算子的本征值；然后将非绝热加热作为微扰项，利用量子力学中发展的微扰论工具，系统分析湿过程对频散关系和波动结构的修正。

\subsection{本文结构}

第 \ref{sec:framework} 节建立理论框架：从三维原始方程出发，通过合理的假设和线性化，导出垂直本征值问题和水平结构方程，阐明解空间的数学结构。
第 \ref{sec:waves} 节讨论干大气波动，复现 Matsuno 理论并分析垂直结构的显式解。
第 \ref{sec:perturbation} 节发展微扰理论，将非绝热加热视为对干波系统的扰动，推导频率和结构的修正。
第 \ref{sec:conclusion} 节进行总结，并简要讨论未来方向。
